% 本文件由 LaTeX 排版 Agent 根据有效 translation.json 自动生成。
% author=Gregory of Nyssa; work_id=2902; title=答欧诺弥乌斯第二书

\chapter{答欧诺弥乌斯第二书}

% structure: p0001 source=h1 target=omitted reason=page-title-already-rendered-by-parent

% structure: p0002 source=h2 target=section reason=relative-heading-rank; base=h2

\section{引言}

为理解以下之书，重要者在于确定心灵中何项官能乃\OriginalTerm{概念}{\GreekText{Ἐπίνοια}}。\Person{欧诺弥乌斯}，\Person{额我略}说，\Quote{对一词作了庄严滑稽的曲解}。他将其力量降至最低水平，仅使之成为\Quote{对反乎自然之物的幻想}，或收缩或延伸自然的界限，或将异质观念拼凑在一起。他举出巨像、侏儒、半人马作为此种心智运作的结果。\Quote{幻想}或\Quote{概念}，由此将代表欧诺弥乌斯对此的看法。但\Person{额我略}将每一门技艺与每一门科学归因于这种官能的运作。\Quote{按我的说法，它是我们借以发现未知之物的方法，通过就所研究之物的\Emph{感知}之周边及后续内容而继续发现。}他一面举出本体论（！）、算术、几何，另一面举出农业、导航、钟表作为其结果。\Quote{任何人若判断此官能比我们所具备的任何其他活动都更为宝贵，则不会大错特错。}\Quote{归纳}几乎可代表此看法。但\Person{额我略}并不否认\Quote{骗人的奇事也由它编造而成}。借此，\Quote{一个表演者可以用喷火怪物、被巨蛇缠绕的人等娱悦观众。}他称之为一种发明性的官能。因此它必比演绎或归纳的推理更为自发，同时比幻想或想象更为可靠。\par

这一论述可由\Person{圣若望·达玛森}在其\WorkTitle{辩证法}（第65章）中对\OriginalTerm{概念}{\GreekText{Ἐπίνοια}}的论述加以说明：\Quote{它有两种。第一种是分析并阐明未经分割总体事物的官能（\OriginalTerm{整体之见}{\GreekText{ὁ} \GreekText{λοσχερῆ}}）：由此一个简单的现象在思辨中变得复杂，例如人成为由灵魂与身体组成的复合体。第二种则藉感知与幻想的联合，从实在中产生虚构，即，将整体分为部分，并将任意选择的部分组合成新的整体，例如半人马、塞壬。}分析（科学的）描述前者；幻想描述后者。\Person{巴西略}与\Person{额我略}思考的是前者，\Person{欧诺弥乌斯}思考的是后者；但双方仍使用同一表述。\par

如果有一个词能涵盖全部含义，那似乎是\Quote{概念}。这个词无论如何，就其外在形式和内在意图，都与感知的关系严格类似，正如\OriginalTerm{概念}{'\GreekText{Επίνοια}}相对于\par

\SourceRecord{Translated by M. Day.；Nicene and Post-Nicene Fathers, Second Series；Vol. 5.；Edited by Philip Schaff and Henry Wace.；Buffalo, NY: Christian Literature Publishing Co.,；1893.；<http://www.newadvent.org/fathers/2902.htm>.}
